\section{Results} \label{sec:results}

\subsection{Financial Proficiency, Savings and Consumption Index, and Behavioral Outcomes} \label{subsec:literacy}

\autoref{tab:ITT} shows that the intervention has a positive and statistically significant impact of 0.07 of a standard deviation on students' financial proficiency, with the effect driven by the middle school grades. Larger impacts on students in more advanced grades are consistent with the hypothesis that financial education programs are more effective at increasing the financial proficiency of students who are more likely to make financial decisions in daily life \citep{Bruhn2016, frisancho2023school, zia2023financial}. School-based programs at large raise financial knowledge by 0.25 of a standard deviation on average \citep{kaiser2020financial}. Our estimate sits at the lower end of the wide range documented for elementary and middle school students, which runs from no detectable effect in Ghanaian primary and junior high schools \citep{berry2018impact} to 0.25 of a standard deviation among American fourth and fifth graders \citep{batty2020experiential} and 0.46 among Belgian middle school students \citep{de2021effect}, although comparisons across programs are complicated by differences in intensity, content, delivery, and setting.

\begin{table}[h!]
\centering
\begin{threeparttable}
\fontsize{9}{10}\selectfont
\caption{ITT estimates for financial proficiency and savings and consumption index}\label{tab:ITT}
\begin{tabular*}{\textwidth}{@{\extracolsep{\fill}}l*{7}{c}}
\toprule
 & (1) & (2) & (3) & (4) & (5) & (6) & (7) \\
 & Pooled & Elementary & Middle school & Third grade & Fifth grade & Seventh grade & Ninth grade \\
\midrule
\tablefragment{../Datawork/Output/Tables/Table3.tex}
\bottomrule
\end{tabular*}
\begin{tablenotes}
\fontsize{9}{9}\selectfont
\item \textit{Notes}: Clustered standard errors at the school level in parentheses. RW denotes Romano-Wolf stepdown p-values (1{,}000 bootstrap replications resampling school clusters within randomization strata), adjusting for the joint testing of the 19 regressions reported in this table (every outcome in every subsample) as a single family, and RI denotes randomization inference (1{,}000 repetitions); both are reported below the standard error. Stars denote significance levels based on the Romano-Wolf adjusted p-value: * 10\%, ** 5\%, *** 1\%. Financial proficiency is computed by CAEd from students' answers to a standardized test scored with Item Response Theory. The consumption index is the first principal component of students' answers to seven agreement-scale statements, and the savings index the first principal component of eight such statements; the statements and their weights are reported in \autoref{tab:cs_index}. All three indices are normalized to standard-deviation units, so treatment effects are read in standard deviations. The questionnaire on hypothetical consumption and saving attitudes was not administered in the third grade. For these two outcomes, column (4) is therefore empty and column (2) reports the same estimates as column (5), the fifth grade.
\item \textit{Source}: Authors' estimates based on the pilot questionnaires applied by CAEd to students at the end of the 2015 school year.
\end{tablenotes}
\end{threeparttable}
\end{table}

Our estimates also point to positive impacts on consumption and saving indices, the ones measuring students' hypothetical attitudes, particularly among fifth graders (\autoref{tab:ITT}). Students in treated classes report more prudent attitudes: they are less likely to endorse statements tolerating debt and impulsive spending, and more likely to endorse the importance of saving. Among fifth graders, the program raised the consumption index by 0.12 of a standard deviation and the savings index by 0.09 (column 5). Among middle school students (column 3), the effect size was smaller for both indices, and statistically significant for consumption index only, suggesting that the program did less to shift the reported savings attitudes of older students.

The larger point estimates among fifth graders might be associated with a few factors. Younger students might be less overconfident than teenagers, making them more receptive to the program's message. In addition, it might be easier to shift the reported attitudes of younger students, whose financial habits are not consolidated yet.  

\newpage
To investigate whether the intervention had distributional impacts on financial proficiency, we estimate unconditional quantile intention-to-treat (UQITT) effects, as described in \autoref{sec:identification}. \autoref{fig:te_firstcycle} displays the UQITT estimates at nineteen quantiles of the financial proficiency distribution, for the pooled sample and separately for elementary and middle school students. Among elementary school students, the estimates become positive and statistically significant from the 45th percentile onward indicating that the program impacted the financial proficiency of a subset of elementary school students. Among middle school students, the UQITT is slightly higher at lower quantiles, but not statistically different from the ITT estimates, suggesting a more homogeneous effect across the financial proficiency distribution in higher grades. 

\begin{figure}[htbp]
		\caption{Unconditional Quantile ITT estimates for financial proficiency - 95\% CI}
		\centering
		\includegraphics[width=\textwidth]{Figure2.pdf}
	    \label{fig:te_firstcycle}
     
     \raggedright
        \fontsize{9}{10} \selectfont

     {\textit{Note}: Financial proficiency is computed by CAEd from students' answers to a standardized test scored with Item Response Theory. The index is normalized to standard-deviation units, so treatment effects are read in standard deviations. We first run this financial proficiency score on the treatment dummy, the six strata dummies, and a constant. We then take the residuals from this estimation and run a simultaneous quantile regression for each quantile shown in the Figure, with the dependent variable being the residual from the first regression and the treatment dummy as the independent variable. 1000 bootstrap replications. The three panels report the pooled sample, elementary school students (third and fifth grades), and middle school students (seventh and ninth grades).}

     {\textit{Source}: Authors' estimates based on the pilot financial-literacy test applied by CAEd to students at the end of the 2015 school year.}
\end{figure}

\newpage

We further explore whether the program affected students' behavioral outcomes. These are self-reported measures, but more objective than the attitudional indeces. \autoref{tab:ITT_beha} shows that fifth graders are 5.7 pp (11.6 percent over the control mean) more likely to talk with their parents about money or family expenses, and 5.6 pp (28.4 percent) more likely to talk about money with their friends. Middle school students are also more likely to talk about money with their friends, but the effect size is small and statistically significant at 10 percent. 

Finally, we test whether the gains in financial proficiency and the more prudent attitudes toward saving and consumption, which might indicate that students became more forward-looking, translated into more investments in human capital. If that were the case, we would hypothesize lower repetition and dropout rates and a higher proficiency in Portuguese and/or math.

We look at grade repetition and dropout between 2015 and 2018; at \textit{Prova Brasil} Portuguese and math scores, measured in the pilot year for the fifth- and ninth-grade cohorts and two years later for the third- and seventh-grade cohorts; and, for the ninth-grade cohort, at high-school completion and \textit{ENEM} outcomes in 2018.

\autoref{tab:long-lasting} shows no systematic evidence that the program affected these outcomes. These outcomes are, however, observed over a relatively short window, one to three years after a low-intensity intervention, so changes that take longer to materialize, especially at these young ages, may fall outside what our data can capture.
 
\begin{table}[ht]
\centering
\begin{threeparttable}
\fontsize{9}{10}\selectfont
\caption{ITT estimates for human capital accumulation outcomes}\label{tab:long-lasting}
\begin{tabular*}{\textwidth}{@{\extracolsep{\fill}}l cccc}
\toprule
 & Third grade & Fifth grade & Seventh grade & Ninth grade \\
 & (1) & (2) & (3) & (4) \\
\midrule
\multicolumn{5}{l}{\textit{Panel A. Grade progression (School Census, cumulative 2015--2018)}} \\
{\bf Repetition} & 0.017 & 0.029** & -0.008 & 0.018 \\
  & (0.017) & (0.014) & (0.016) & (0.018) \\
{\bf Dropout} & 0.005 & 0.011 & 0.002 & 0.020 \\
  & (0.006) & (0.010) & (0.012) & (0.019) \\
\midrule
\multicolumn{5}{l}{\textit{Panel B. Test scores (\textit{Prova Brasil})}} \\
\textit{Assessed in} & \textit{2017} & \textit{2015} & \textit{2017} & \textit{2015} \\
{\bf Portuguese} & -0.029 & -0.017 & -0.005 & 0.028 \\
  & (0.047) & (0.046) & (0.034) & (0.036) \\
{\bf Math} & -0.066 & -0.033 & -0.009 & -0.020 \\
  & (0.053) & (0.049) & (0.038) & (0.039) \\
\midrule
\multicolumn{5}{l}{\textit{Panel C. End of high school (ninth-grade cohort assessed in 2018)}} \\
{\bf Graduated from high school} &  &  &  & -0.019 \\
  &  &  &  & (0.023) \\
{\bf Participated in ENEM} &  &  &  & -0.012 \\
  &  &  &  & (0.023) \\
{\bf Portuguese ENEM} &  &  &  & -0.025 \\
  &  &  &  & (0.031) \\
{\bf Math ENEM} &  &  &  & 0.002 \\
  &  &  &  & (0.045) \\
\midrule
Obs. & 4650 & 4755 & 3137 & 2851 \\
\bottomrule
\end{tabular*}
\begin{tablenotes}
\fontsize{9}{9}\selectfont
\item \textit{Notes}: Clustered standard errors at the school level in parentheses. Stars denote significance levels based on conventional p-values: * 10\%, ** 5\%, *** 1\%. The estimates were produced inside INEP's secure facility, where the bootstrap-based Romano-Wolf were not run; since the estimates are null, adjusting for multiple hypothesis testing would only reinforce the conclusions. Each dependent variable is regressed on treatment assignment, with one specification per cohort. Columns (1) and (2) include strata fixed effects only, since the third- and fifth-grade cohorts have no pre-pilot test score. Columns (3) and (4) add the student's pre-pilot fifth-grade \textit{Prova Brasil} performance, taken in 2013 by the seventh-grade cohort and in 2011 by the ninth-grade cohort, which controls for the pre-treatment imbalance documented in \autoref{tab:balance_adm2}.
\item \textit{Timing}: repetition and dropout (Panel A) are cumulative over 2015--2018. In Panel B, the ``Assessed in'' row gives the year each cohort sits \textit{Prova Brasil}: the fifth- and ninth-grade cohorts are assessed in the pilot year (2015, contemporaneous), while the third- and seventh-grade cohorts are assessed two years later (2017), when they reach the next assessment grade (see \autoref{tab:inep_data}). Panel C outcomes refer to 2018 and exist only for the ninth-grade cohort.
\item \textit{Source}: Authors' estimates based on the pilot roster merged with INEP's identified student-level administrative records (School Census, \textit{Prova Brasil}, and ENEM). INEP performed the merge; the authors ran all estimates inside INEP's secure facility.
\end{tablenotes}
\end{threeparttable}
\end{table}


Our estimates show that the program had larger impacts on financial proficiency in later grades and larger effects on attitudes and behavior in earlier grades. The results then suggest that the association between financial proficiency and behavioral change might not be straightforward. Individuals might have formal knowledge across different subjects, such as financial proficiency, but the application of that knowledge in real life can be hindered by the institutional environment they are embedded in \citep{carpena2020causal}. To formally test the hypothesis that learning mediates change in behavior, we follow \cite{imai2011unpacking} and
decompose the treatment effect estimates into the average causal mediation effect (ACME) and
the average causal direct effect (ADE).\footnote {See the online appendix for a formal discussion of Causal Mediation Effects.} The goal of this exercise is to assess the share of the change in behavioral outcomes that could be attributed to improvements in financial proficiency (ACME) and the share that comes directly from the program (ADE).\footnote{The decomposition of ITT effects on ACME and ADE relies on the sequential ignorability assumption. The first part of the assumption assumes that, conditioned on pre-treatment confounders, the treatment assignment is orthogonal to the potential outcomes and the mediator outcome (financial proficiency). The second part of the assumption implies that the observed mediator is ignorable given the actual treatment status and pre-treatment confounders \citep{imai2011unpacking}. This second part of the assumption is the strongest, even in the presence of random assignment. It posits that, conditional on the actual treatment (not the treatment assignment) and a vector of pre-treatment confounders, the outcome is independent of the mediator. Note that in our case, this implies the assumption that attitudinal and behavioral outcomes are orthogonal to financial proficiency, given the actual treatment and after controlling for a vector of school characteristics. Since we did not design the experiment to obtain clean estimates of the ACME and ADE, we interpret our results as suggestive evidence.}

\begin{table}[h!]
\centering
\begin{threeparttable}
\fontsize{9}{10}\selectfont
\caption{ITT estimates for behavioral outcomes}\label{tab:ITT_beha}
\begin{tabular*}{\textwidth}{@{\extracolsep{\fill}}l*{3}{c}}
\toprule
 & (1) & (2) & (3) \\
 & Pooled & Elementary & Middle school \\
\midrule
\tablefragment{../Datawork/Output/Tables/Table4.tex}
\bottomrule
\end{tabular*}
\begin{tablenotes}
\fontsize{9}{9}\selectfont
\item \textit{Notes}: Clustered standard errors at the school level in parentheses. RW denotes Romano-Wolf stepdown p-values (1{,}000 bootstrap replications resampling school clusters within randomization strata), adjusting for the joint testing of the 15 regressions reported in this table (every outcome in every subsample) as a single family, and RI denotes randomization inference (1{,}000 repetitions); both are reported below the standard error. Stars denote significance levels based on the Romano-Wolf adjusted p-value: * 10\%, ** 5\%, *** 1\%. Talking about money or expenses with parents equals 1 if the student reports discussing household expenses or talking about money with their parents, and 0 otherwise; the remaining variables equal 1 if the student talks about money with friends, keeps a piggy bank, has access to at least one financial service (bank account, debit card, or credit card), or receives an allowance from their parents, and 0 otherwise. The sample is restricted to fifth-, seventh-, and ninth-grade students.
\item \textit{Source}: Authors' estimates based on the pilot questionnaires applied by CAEd to students at the end of the 2015 school year.
\end{tablenotes}
\end{threeparttable}
\end{table}


According to the estimates in \autoref{tab:acme}, financial proficiency mediates part of the improvement in students' hypothetical attitudes toward consumption: among middle school students, the mediated component accounts for about half of the total effect. For elementary school students, by contrast, every mediated component is statistically indistinguishable from zero, so the effects on attitudes and on talking about money with parents and friends are direct. The same holds for the behavioral outcomes at both school levels, where the mediated components are negligible. In other words, the ADE and ACME estimates suggest that financial education programs affect students' attitudes and behavior through pathways other than the accumulation of formal knowledge. Overall, the estimates do not support the hypothesis that changes in financial proficiency are a necessary condition for changes in individuals' behavioral outcomes, a finding consistent with \citet{carpena2020causal} and \citet{de2021effect}.



\subsection{Robustness Check} 

As discussed in \autoref{sec:data_c}, the student-level merge with INEP's records recovers pre-treatment proficiency scores in reading and math for the seventh- and ninth-grade cohorts, and \autoref{tab:balance_adm2} shows that ninth-grade control students outperformed treated ones in reading before the pilot. Since principals, and not chance, chose the participating classes, this imbalance could bias our estimates. We therefore re-estimate \autoref{eq:Y} for these two cohorts controlling for their fifth-grade proficiency scores.\footnote{These are the only two cohorts with a pre-treatment performance score: seventh graders sat the fifth-grade test in 2013 and ninth graders in 2011.}

Controlling for past performance leads to a substantial gain in precision. Among ninth graders, the point estimate for financial proficiency rises from 0.107 to 0.149 of a standard deviation; and among seventh graders, it rises slightly from 0.09 to 0.098 (\autoref{tab:add_5prof_control}). The pattern suggests that ignoring the pre-treatment imbalance would underestimate the program's impact on financial proficiency, particularly for ninth graders. Controlling for past performance does not change the estimates for consumption and saving attitudes. 


\subsection{Implementation and IV Results}

As discussed in \autoref{sec:implementation}, we found evidence of contamination, especially in the control group. Because exposure to financial education content varied across classes within schools in both arms, the ITT compares assignment rather than exposure. To recover the effect of actual exposure, we use the random assignment as an instrument for exposure, following \cite{frisancho2023school}.

We measure exposure using two definitions, regardless of the arm to which the school was assigned. The conservative one is defined at the class level: it sets $Z=1$ when a teacher reports delivering financial education classes.\footnote{If the teacher answered the question ``\textit{How many financial education classes did you teach?}'', with answer categories running from one to eight or more, or the question ``\textit{What share of the student textbook did you cover?}'', with categories running from below 40\% to above 80\%.} For the classes whose teachers did not answer the teacher questionnaire, we set $Z=1$ if more than half of the students from those classes report having had such classes. The broad one is defined at the student level: it sets $Z=1$ when a teacher of the student's class reports delivering the financial education classes, as in the conservative definition, or when the student themselves reports having had them. Under the conservative definition, 97.8\% of students in treated classes and 11.1\% in control classes are classified as exposed; under the broad definition, 97.3\% and 21.5\%, respectively.

We estimate the local average treatment effect (LATE) by two-stage least squares. The first stage is
\begin{align}
     Z_{ism} = \pi_0 + \pi_1 T_{ism} + \sum_{n=1}^{6} \gamma_n + \nu_{ism}
     \label{eq:Y2}
\end{align}
and the second stage is
\begin{align}
    Y_{ism} = \lambda_0 + \lambda_1 \widehat{Z}_{ism} + \sum_{n=1}^{6} \gamma_n + \omega_{ism}
    \label{eq:Y3}
\end{align}
\noindent
where $Z$ is the exposure of student $i$'s class under the conservative definition and the student's own exposure under the broad one, $T$ is the random assignment of \autoref{eq:Y}, and the remaining notation follows that equation. Standard errors are clustered at the school level, the randomization unit. The parameter of interest, $\lambda_1$, is the LATE: the average effect of exposure among compliers, that is, the classes (under the conservative definition) or students (under the broad definition) that were exposed to the program because their school was assigned to the treatment group, and would not have been exposed otherwise.

\autoref{tab:iv} reports the LATE estimates next to the ITT under both definitions of exposure.\footnote{The first-stage estimates behind each LATE are reported in \autoref{tab:first_stage}.} For financial proficiency the differences are moderate: among middle school students, the effect rises from 0.099 of a standard deviation to 0.112 under the conservative definition and 0.135 under the broad one. For the attitude indices in grade 5, the correction is substantial: under the broad definition, the effect on the consumption index rises from 0.122 to 0.202 of a standard deviation, and on the saving index from 0.087 to 0.143. Exposure to the program therefore had modest effects on financial proficiency even among the classes that delivered the content, while its effects on the reported attitudes of younger students are sizeable once imperfect compliance is taken into account; since the broad definition relies on students' own reports, these larger estimates are best read as upper bounds.

\vspace{10 pt}

\begin{table}[h!]
\centering
\begin{threeparttable}
\fontsize{9}{10}\selectfont
\caption{Intention-to-treat (ITT) and local average treatment effect (LATE) estimates}\label{tab:iv}
\setlength{\tabcolsep}{2.5pt}
\begin{tabular*}{\textwidth}{@{\extracolsep{\fill}}l*{9}{c}}
\toprule
 & \multicolumn{3}{c}{Pooled} & \multicolumn{3}{c}{Elementary} & \multicolumn{3}{c}{Middle school} \\
\cmidrule(lr){2-4}\cmidrule(lr){5-7}\cmidrule(lr){8-10}
 & ITT & \shortstack{LATE\\ (i)} & \shortstack{LATE\\ (ii)} & ITT & \shortstack{LATE\\ (i)} & \shortstack{LATE\\ (ii)} & ITT & \shortstack{LATE\\ (i)} & \shortstack{LATE\\ (ii)} \\
\midrule
\tablefragment{../Datawork/Output/Tables/Table6.tex}
\bottomrule
\end{tabular*}
\begin{tablenotes}
\fontsize{9}{9}\selectfont
\item \textit{Notes}: Clustered standard errors at the school level in parentheses. RW denotes Romano-Wolf stepdown p-values (1{,}000 bootstrap replications resampling school clusters within randomization strata), adjusting for the joint testing of the 27 regressions reported in this table (every outcome, subsample, and estimator) as a single family, reported below the standard error. Stars denote significance levels based on the Romano-Wolf adjusted p-value: * 10\%, ** 5\%, *** 1\%. ITT is the reduced-form effect of random assignment; LATE instruments actual program exposure (\(Z\)) with random assignment. Exposure is measured under two definitions: (i) the conservative one, defined at the class level, which requires a teacher of the class to report having taught the financial education classes or covered part of the textbook, and, for the classes whose teachers did not answer their pilot questionnaire, that more than half of the students report having had such classes; and (ii) the broad one, defined at the student level, which sets \(Z=1\) when a teacher of the class reports delivering the classes or the student themselves reports having had them. Both are defined in the text. All three indices are normalized to standard-deviation units. The consumption and savings indices are unavailable for third graders, who did not answer these items. Therefore, elementary estimates represent the fifth grade, while middle school estimates represent the seventh and ninth grades. 
\item \textit{Source}: Authors' estimates based on the pilot questionnaires applied by CAEd to students at the end of the 2015 school year.
\end{tablenotes}
\end{threeparttable}
\end{table}
